\section{Five-Way Equivalence}\label{five-way-equivalence}

This section establishes the central result: five independent characterizations of the same architectural property (DOF = 1 / SSOT) derived from five distinct first-principles frameworks. All proofs are machine-checked in Lean 4.

\subsection{Framework 1: Engineering Optimization (Leverage)}

\begin{formal}
\textbf{Theorem 5.1 (Maximum Leverage).} For any architecture $A$ with $\text{Cap}(A) > 0$, $A$ achieves maximum leverage among architectures with equal capabilities if and only if $\text{DOF}(A) = 1$.

Formally: $L(A) = \max_{A': \text{Cap}(A')=\text{Cap}(A)} L(A') \iff \text{DOF}(A) = 1$.
\end{formal}

\begin{informal}
An architecture has the highest possible capability-to-DOF ratio exactly when it has a single degree of freedom.
\end{informal}

\textbf{Proof.} 
\begin{itemize}
\item (Forward) If $\text{DOF}(A) = 1$, then $L(A) = |\text{Cap}(A)|/1 = |\text{Cap}(A)|$. Any $A'$ with same capabilities has $\text{DOF}(A') \geq 1$, so $L(A') \leq |\text{Cap}(A)| = L(A)$.
\item (Backward) If $A$ has maximum leverage but $\text{DOF}(A) > 1$, construct $A'$ with $\text{DOF}(A') = 1$ and same capabilities. Then $L(A') = |\text{Cap}(A)| > |\text{Cap}(A)|/\text{DOF}(A) = L(A)$, contradiction.
\end{itemize}

\leanmeta{\LH{L28}, \LH{L31}, \LH{L48}}

\subsection{Framework 2: Architectural Epistemic Coherence (Ssot)}

\begin{formal}
\textbf{Theorem 5.2 (Coherence).} An architecture satisfies the Single Source of Truth property if and only if $\text{DOF}(A) = 1$.

Formally: $\text{SSOT}(A) \iff \text{DOF}(A) = 1$.
\end{formal}

\begin{informal}
Epistemic coherence (one source of truth per structural fact) is equivalent to having a single degree of freedom.
\end{informal}

\leanmeta{Proved in \nolinkurl{Ssot} (\LH{ORA1})}

This is the DOF = 1 characterization from the Ssot coherence theorem.

\subsection{Framework 3: Information-Theoretic Structural Rank (DecisionQuotient)}

\begin{formal}
\textbf{Theorem 5.3 (DOF-Structural Rank Isomorphism).} For any architecture $A$, let $\text{canonicalDP}(n)$ be the canonical decision problem with $n$ boolean coordinates. Then:
\[
\text{srank}(\text{canonicalDP}(\text{DOF}(A))) = \text{DOF}(A)
\]
\end{formal}

\begin{informal}
The structural rank (number of relevant decision coordinates) equals the degrees of freedom. The canonical encoding uses: states as $n$ boolean coordinates, actions as either querying a coordinate or falling back, and utilities that reward correct coordinate identification.
\end{informal}

The canonical encoding is:
\begin{itemize}
\item States: $\text{Fin } n \to \text{Bool}$ ($n$ binary coordinates)
\item Actions: $\text{Fin } n \oplus \text{Unit}$ (query coordinate $i$, or fallback)
\item Utility: $u(\text{inl } i, s) = 2$ if $s(i) = \text{true}$, else $0$; $u(\text{inr}\ (), s) = 1$
\end{itemize}

\leanmeta{\LH{L43}, \LH{L46}, \LH{L51}}

\begin{formal}
\textbf{Corollary 5.4.} DOF $= 1$ if and only if srank $= 1$.

Formally: $\text{DOF}(A) = 1 \iff \text{srank}(\text{canonicalDP}(\text{DOF}(A))) = 1$.
\end{formal}

\begin{informal}
An architecture has a single degree of freedom exactly when its canonical decision problem has minimal structural rank.
\end{informal}

\leanmeta{Immediate from Theorem 5.3}

\subsection{Framework 4: Computational Complexity (Tractability)}

\begin{formal}
\textbf{Theorem 5.5 (Tractability Boundary).} For the canonical decision problem with $n$ coordinates:
\begin{enumerate}
\item If $n = 1$ (srank $= 1$), sufficiency checking is in P (polynomial time).
\item If $n > 1$ (srank $> 1$), sufficiency checking is coNP-hard.
\end{enumerate}
\end{formal}

\begin{informal}
When the structural rank is 1, determining whether current information suffices for an optimal decision is computationally tractable. When the structural rank exceeds 1, the same problem is computationally intractable (coNP-hard).
\end{informal}

\leanmeta{Proved in \texttt{DecisionQuotient} (\LH{L56}, \LH{L53}), imported via \nolinkurl{Leverage/BridgeToDQ.lean}}

This connects the architectural property (DOF = 1) to computational feasibility of decision-making.

\begin{formal}
\textbf{Corollary 5.6.} DOF $= 1$ (SSOT) is the unique architecture class for which sufficiency checking is tractable.
\end{formal}

\begin{informal}
Among architectures, only those with a single degree of freedom admit tractable sufficiency checking.
\end{informal}

\leanmeta{Combine Corollary 5.4 with Theorem 5.5; \LH{L47}}

\subsection{Framework 5: Thermodynamic Selection (Statistical Physics)}

\begin{formal}
\textbf{Theorem 5.7 (Thermodynamic Selection).} Let $M$ be a thermodynamic model with Landauer calibration. For any architecture $A$ with DOF $> 1$:
\begin{enumerate}
\item There exist decision instances where sufficiency checking requires $\Omega(2^{\text{DOF}})$ logical operations.
\item Under the physical axioms of \cite{paper4_decision_quotient} (\LH{LP38}): Landauer (\LH{EP1}), nontrivial state space (\LH{EP4}), finite signal speed (\LH{LP43}), finite universe energy (\LH{LP44}) — P~$\neq$~NP is proved, hence P~$\neq$~coNP; no polynomial-time sufficiency certification exists.
\item \textbf{(Unconditional)} Each sufficiency-check cycle incurs mandatory positive energy cost $\geq j_{\mathcal{M}} \cdot \mathrm{srank}(A) > j_{\mathcal{M}}$, proved from Landauer calibration and the DOF~$=$~srank identity alone, with no complexity hypothesis.
\end{enumerate}
For DOF $= 1$, the energy lower bound is $j_{\mathcal{M}}$ (the Landauer minimum).
\end{formal}

\begin{informal}
Architectures with more than one degree of freedom necessarily incur thermodynamic costs per decision cycle. Point~3 holds unconditionally: the energy bound follows from Landauer and the DOF~$=$~srank identity (\LH{L43}). Point~2 holds under the physical axioms of \cite{paper4_decision_quotient} (\LH{LP38}), which prove P~$\neq$~NP, hence P~$\neq$~coNP.
\end{informal}

Physical assumptions (point~3 only):
\begin{itemize}
\item Positive Landauer constant ($j_{\mathcal{M}} > 0$) — cf. Theorem \ref{cor:leverage-energy} (\LH{L3}, \LH{L6})
\item Landauer calibration ($k_B T \ln 2$ joules per bit)
\end{itemize}

\leanmeta{\LH{L49}, \LHrng{L}{54}{55}}

Physical edit-energy floor from Section 4.3 is a special case: DOF controls minimum edit-energy, and higher leverage implies lower energy in a fixed capability class (Theorem \ref{thm:physical-energy-floor}, \LH{L1}, \LH{L5}).

\subsection{The Five-Way Equivalence}

\begin{formal}
\textbf{Theorem 5.8 (Five-Way Equivalence).} For any architecture $A$ with $\text{Cap}(A) > 0$, the following are equivalent:
\[
\text{DOF}(A) = 1 \iff \text{max leverage} \iff \text{SSOT}(A) \iff \text{srank}(A) = 1 \iff \text{tractable sufficiency} \iff \text{minimum thermodynamic cost}
\]
\end{formal}

\begin{informal}
Five independent scientific frameworks—engineering optimization, epistemic coherence, information geometry, computational complexity, and statistical physics—all select the same architectural property: having exactly one degree of freedom.
\end{informal}

\begin{center}
\begin{tabular}{ll}
\hline
\textbf{Domain} & \textbf{Characterization of DOF $= 1$} \\
\hline
Engineering Optimization & Maximum leverage \\
Architectural Epistemic & Single Source of Truth \\
Information Geometry & Structural rank $= 1$ \\
Computational Complexity & Tractable sufficiency checking \\
Statistical Physics & Minimum thermodynamic cost per cycle \\
\hline
\end{tabular}
\end{center}

\textbf{Proof.} The equivalence follows from Theorems 5.1-5.7:
\begin{enumerate}
\item $\text{DOF} = 1 \iff \text{max leverage}$: Theorem 5.1 (\LH{L31}, \LH{L28})
\item $\text{DOF} = 1 \iff \text{SSOT}$: Theorem 5.2 \cite{paper2_ssot}
\item $\text{DOF} = \text{srank}$: Theorem 5.3, so $\text{DOF} = 1 \iff \text{srank} = 1$
\item $\text{srank} = 1 \iff \text{tractable sufficiency}$: Theorem 5.5 \cite{paper4_decision_quotient}. Forward direction (srank~$=1 \to$ polynomial) is unconditional. Backward direction (tractable $\to$ srank~$=1$) holds under P~$\neq$~coNP, which follows from the physical axioms of \cite{paper4_decision_quotient} (\LH{LP38}).
\item $\text{DOF} > 1 \iff \text{above-minimum thermodynamic cost}$: Theorem 5.7, so $\text{DOF} = 1 \iff \text{minimum thermodynamic cost}$
\end{enumerate}

The transitivity of logical equivalence gives the five-way equivalence.

\leanmeta{\LHrng{L}{43}{44}, \LH{L47}, \LH{L55}, \LH{ORA1}}

\begin{formal}
\textbf{Corollary 5.9 (Convergence).} Within the \nolinkurl{DecisionQuotient} canonical decision problem framework, the boolean-coordinate encoding (states as $n$ boolean coordinates, actions as coordinate queries or fallback, utilities rewarding correct identification) is the unique decision structure with $n = 1$ that simultaneously satisfies all five characterizations.
\end{formal}

\begin{informal}
In the canonical decision problem representation, there is exactly one $n=1$ structure, and all five characterizations select it. This is a consequence of the equivalences, not an additional theorem. Note: there are many architectures with DOF~$=1$ (monolith, single DB table, single endpoint); the uniqueness claim is scoped to the canonical decision problem encoding.
\end{informal}

\textbf{Proof.} Each characterization is equivalent to DOF~$= 1$ (Theorems 5.1--5.7). Within the canonical encoding, DOF~$= 1$ fixes $n = 1$, and $\mathrm{Fin}\ 1 \to \mathrm{Bool}$ has a unique element. All five requirements therefore select the same structure.

\leanmeta{All five requirements force $n=1$ by the above theorems; uniqueness within the canonical encoding follows from $|\mathrm{Fin}\ 1 \to \mathrm{Bool}| = 2$}

\subsection{Formalization}

The following Lean 4 proofs establish the connections:

\begin{itemize}
\item \nolinkurl{Leverage/Foundations.lean}:
  \begin{itemize}
  \item \LH{L31}, \LH{L28} (\LH{L50}): DOF $= 1 \to$ max leverage
  \item \LH{L48}: max leverage $\to$ DOF $= 1$
  \item \LH{L44}: biconditional
  \end{itemize}
\item \nolinkurl{Leverage/BridgeToDQ.lean}:
  \begin{itemize}
  \item \LH{L43}: DOF $=$ srank
  \item \LH{L51}: DOF $= 1 \to$ srank $= 1$
  \item \LH{L46}: DOF $> 1 \to$ srank $> 1$
  \item \LH{L54}: DOF $> 1 \to$ mandatory energy
  \item \LH{L47}: max leverage $\to$ tractable
  \end{itemize}
\item \nolinkurl{DecisionQuotient/Tractability/StructuralRank.lean}: srank $= 1 \to$ P (\LH{L56}), srank $> 1 \to$ coNP-hard (\LH{L53})
\item \nolinkurl{DecisionQuotient/ThermodynamicLift.lean}: energy lower bounds under Landauer calibration (\LH{L49})
\end{itemize}

The cross-module dependency chain is live: \texttt{Leverage} $\to$ \texttt{DecisionQuotient} $\to$ Mathlib.

\subsection{Relation to Prior Sections}

This section subsumes and extends Section 4's results:
\begin{itemize}
\item Theorem 4.1 (Leverage-Error Tradeoff, \LH{L29}) is a consequence of the equivalence between leverage and DOF
\item Theorem 4.2 (Modification Complexity Gap, \LH{L30}) follows from DOF $= 1$ minimizing modification cost
\item Theorem 4.3 (Optimal Architecture, \LH{L38}, \LH{L34}) is strengthened: the optimal architecture is now characterized by five independent properties
\item Theorem 4.4 (Metaprogramming Dominance, \LH{L36}, \LH{L35}) is a special case: unbounded derivations from a single source achieve DOF $= 1$
\end{itemize}

\subsection{England Replication Inequality (Proved)}

All previously open conjectures are now proved. The England Replication Inequality is mechanized in \nolinkurl{Leverage/BridgeToDQ.lean} (\LH{L45}).

\begin{theorem}[England Replication Inequality]\label{thm:england}
Let $\Delta S_{\min}(\text{srank}) = \text{srank} \cdot k_B \ln 2$ be the minimal entropy production under Landauer calibration. For a single-source architecture ($\text{srank} = 1$) and a $k$-copy replication architecture ($\text{srank} = k$):
\[
\Delta S_{\min}(1) + k_B \ln k \leq \Delta S_{\min}(k)
\]
equivalently, $\Delta S_{\min}(k) - \Delta S_{\min}(1) \geq k_B \ln k$.
\end{theorem}

\leanmeta{\LH{L45}}

\begin{proof}
The gap is $(k-1) \cdot k_B \ln 2$. Since $k \leq 2^{k-1}$ (\LH{L52}), taking logs gives $\ln k \leq (k-1) \ln 2$, so the gap is $\geq k_B \ln k$.
\end{proof}

\textbf{Modeling note.} $\Delta S_{\min}$ is a definition within the model: the exact Landauer entropy cost per cycle, not a lower bound on arbitrary implementations. The ``min'' refers to physical optimality outside the formalism.
